\newif\ifnotesonly
\notesonlyfalse % \notesonlytrue

\documentclass[aspectratio=169, 10pt, xcolor=svgnames]{beamer}

\ifnotesonly
  \setbeameroption{show only notes}
\fi

% Orchid-Thistle Thema
\useoutertheme[subsection=false]{miniframes}
\useinnertheme{rounded}
\setbeamertemplate{navigation symbols}{}

\definecolor{PlumDark}{HTML}{4A2050}
\setbeamercolor{normal text}{fg=PlumDark, bg=white}
\setbeamercolor{structure}{fg=DarkOrchid}
\setbeamercolor{alerted text}{fg=MediumVioletRed}
\setbeamercolor{section in head/foot}{fg=PlumDark, bg=Thistle}
\setbeamercolor{frametitle}{fg=PlumDark, bg=Lavender}
\setbeamercolor{title}{fg=PlumDark, bg=Thistle}
\setbeamercolor{subtitle}{fg=DarkOrchid}
\setbeamercolor{block title}{fg=white, bg=Orchid}
\setbeamercolor{block body}{fg=PlumDark, bg=Lavender}
\setbeamercolor{block title alerted}{fg=white, bg=MediumOrchid}
\setbeamercolor{block body alerted}{fg=PlumDark, bg=Thistle!40}
\setbeamercolor{block title example}{fg=white, bg=Plum!90!black}
\setbeamercolor{block body example}{fg=PlumDark, bg=Thistle!25}
\setbeamercolor{footnote}{fg=DarkOrchid}
\setbeamercolor{footnote mark}{fg=DarkOrchid}
\setbeamercolor{page number in head/foot}{fg=DarkOrchid}
\setbeamercolor{bibliography entry author}{fg=PlumDark}
\setbeamercolor{bibliography entry note}{fg=DarkOrchid}
\setbeamerfont{footnote}{size=\tiny}
\setbeamerfont{frametitle}{series=\bfseries}
\setbeamertemplate{frametitle}{%
  \nointerlineskip
  \begin{beamercolorbox}[wd=\paperwidth, ht=2.6ex, dp=1.2ex, leftskip=0.8em]{frametitle}
    \usebeamerfont{frametitle}\insertframetitle
  \end{beamercolorbox}}
% page numbers from the second slide (title slide — [plain])
\setbeamertemplate{footline}{%
  \hfill\usebeamercolor[fg]{page number in head/foot}\usebeamerfont{footnote}%
  \insertframenumber\,/\,\inserttotalframenumber\hspace{1.2em}\vspace{1.0ex}}
\setbeamertemplate{bibliography item}[text]
% notes page: a thumbnail of the slide at the top, with the narrative text below
\setbeamerfont{note page}{size=\small}
\setbeamertemplate{note page}{%
  \vspace*{0.4em}%
  \fbox{\insertslideintonotes{0.30}}\hspace{1em}%
  \parbox[b]{0.55\paperwidth}{\usebeamercolor[fg]{structure}\bfseries\insertframetitle\par
    \normalfont\color{PlumDark}\footnotesize слайд \insertframenumber}\par
  \vspace{0.6em}%
  \parbox{\textwidth}{\usebeamerfont{note page}\insertnote}}

% Packages
\usepackage{fontspec}
% \setmainfont{Times New Roman}
% \setsansfont{Tahoma}
% \setmonofont{Roboto Mono}
\IfFontExistsTF{CMU Sans Serif}{\setsansfont{CMU Sans Serif}}{%
  \IfFontExistsTF{DejaVu Sans}{\setsansfont{DejaVu Sans}}{\setsansfont{Arial}}}
\IfFontExistsTF{Roboto Mono}{\setmonofont{Roboto Mono}}{}
\usepackage[english]{babel}
\babelprovide[import, main]{ukrainian}
\usepackage{csquotes}
\usepackage{amsmath, amssymb}
\usepackage{booktabs}
\usepackage{graphicx}
\usepackage{tikz}
\usetikzlibrary{arrows.meta, positioning, calc, fit, backgrounds, shapes.geometric}
\usepackage{pgfplots}
\pgfplotsset{compat=1.18}
\usepackage[backend=biber,style=numeric,sorting=none,autolang=other]{biblatex}

\addbibresource{../references.bib}

\newcommand{\G}{\mathbb{G}}
\newcommand{\Zq}{\mathbb{Z}_q}
\newcommand{\UID}{\mathrm{UID}}
\newcommand{\photo}[2][height=0.32\textheight]{\IfFileExists{screenshots/#2}{\includegraphics[#1]{screenshots/#2}}{}}

\title{Довести все, не розкривши нічого:}
\subtitle{Як BBS-підписи та Signal zkgroup ламають правила стеження}
\author{Юрчук Олексій}
\institute{ФІ-52мн}
\date{13.11.2026}

\begin{document}

%  1 slide

\begin{frame}[plain]
  \titlepage
\end{frame}

%  2. План

\begin{frame}{План доповіді}
  \begin{columns}[c]
    \begin{column}{0.55\textwidth}
      \begin{enumerate}
        \item \textbf{Вступ:} Дві родини анонімних облікових даних
        \item \textbf{Проєкт:} Приватні групи Signal -- що і навіщо доводиться
        \item \textbf{Протокол:} Алгебраїчний MAC $+$ $\Sigma$-протоколи
        \item \textbf{Аналіз:} Припущення, ризики, бенчмарки, альтернатива BBS
        \item \textbf{Висновки} і питання для самоперевірки
      \end{enumerate}
      \textbf{Чому ця тема:} Актуально; відсилка до лекції~3; Відкрита програмна реалізація.
    \end{column}
    \begin{column}{0.40\textwidth}
      \begin{block}{Головне питання}
        Як довести серверу, що я \emph{учасник групи}, не кажучи, \emph{хто я}?
      \end{block}
    \end{column}
  \end{columns}
\end{frame}

%  3. Вступ: дві родини
\section{Вступ}
\begin{frame}{Анонімні облікові дані: дві родини}
  \begin{columns}[T]
    \begin{column}{0.58\textwidth}
      \centering
      \resizebox{\linewidth}{!}{\begin{tikzpicture}[
    font=\normalsize,
    node distance=3mm and 4mm,
    root/.style={draw=DarkOrchid, fill=Thistle, rounded corners=3pt, font=\normalsize\bfseries, align=center},
    fam/.style={draw=DarkOrchid, fill=Thistle!45, rounded corners=3pt, align=center, text width=4.6cm},
    leaf/.style={draw=Plum, fill=white, rounded corners=2pt, align=left, font=\small, text width=4.6cm},
    hl/.style={leaf, draw=DarkOrchid, very thick, fill=Lavender},
    edge/.style={DarkOrchid, thick}
  ]
  \node[root] (R) {Анонімні облікові дані (anonymous credentials)\\[-1pt]\small\mdseries Chaum, 1985};
  \node[fam, below left=6mm and -1.9cm of R]  (PV) {\textbf{Публічно верифіковані}\\[-1pt]
        \small (publicly verifiable): \\ облікові дані $=$ \emph{підпис}};
  \node[fam, below right=6mm and -1.9cm of R] (KV) {\textbf{З верифікацією за ключем}\\[-1pt]
        \small (KVAC): облікові \\ дані $=$ \emph{алгебраїчний MAC}};
  \node[leaf, below=of PV] (pv1) {Підписи CL (strong RSA, 2001)};
  \node[hl,   below=1.2mm of pv1] (pv2) {BBS / BBS+ (спарювання, $q$-SDH);\\ чернетка IETF CFRG, BLS12-381};
  \node[leaf, below=1.2mm of pv2] (pv3) {PS, U-Prove, \dots};
  \node[leaf, below=of KV] (kv1) {CMZ: $\mathsf{MAC}_{\mathsf{GGM}}$, $\mathsf{MAC}_{\mathsf{DDH}}$ (2014)};
  \node[hl,   below=1.2mm of kv1] (kv2) {Chase--Perrin--Zaverucha (2020)\\ $\rightarrow$ Signal \texttt{zkgroup}};
  \node[leaf, below=1.2mm of kv2] (kv3) {BBS-MAC (2016), $\mu$CMZ / $\mu$BBS (2024),\\ Tor Lox, IETF ARC};
  \draw[edge] (R) -- (PV); \draw[edge] (R) -- (KV);
  \draw[edge] (PV) -- (pv1); \draw[edge] (KV) -- (kv1);
  \node[font=\small, align=center, below=3mm of $(pv3.south)!0.5!(kv3.south)$]
       {спільний «двигун»: $\Sigma$-протоколи $+$ Фіат--Шамір};
\end{tikzpicture}
}
    \end{column}
    \begin{column}{0.41\textwidth}
      \small
      \textbf{Ідея}\footnotemark[1]: "Посвідчення"{} від видавця (issuer); показуємо лише потрібне, пред'явлення
      не зв'язуються.\par\smallskip
      \textbf{Аналогія:} Паспорт перевіряє будь-хто; жетон гардероба -- лише той гардероб, що його видав
      (за ключем\footnotemark[2]).\par\smallskip
      \centering\photo[width=\linewidth, height=0.42\textheight, keepaspectratio]{analogy_passport_token.jpg}
    \end{column}
  \end{columns}
  \footnotetext[1]{Chaum, CACM 1985 \cite{Chaum85}}
  \footnotetext[2]{Chase, Meiklejohn, Zaverucha, CCS 2014 \cite{CMZ14}}
\end{frame}

%  4. Проєкт: задача
\section{Проєкт}
\begin{frame}{Задача: приватні групи (private groups)}
  \begin{columns}[T]
    \begin{column}{0.68\textwidth}
      \begin{columns}[T]
        \begin{column}{0.49\linewidth}
          \begin{block}{До редизайну}
            \footnotesize
            \begin{itemize}
              \item Групи існують лише на телефонах -- сервер про неї не знає\footnotemark[1]
              \item Немає канонічного стану
              \item Немає ролей (адмін / учасник)
            \end{itemize}
          \end{block}
        \end{column}
        \begin{column}{0.49\linewidth}
          \begin{block}{Private Group System}
            \footnotesize
            \begin{itemize}
              \item Стан групи на сервері, але \textbf{зашифрований}\footnotemark[2]
              \item сервер перевіряє, що запит від учасника, але не знає, від якого саме
            \end{itemize}
          \end{block}
        \end{column}
      \end{columns}
      \medskip
      \begin{alertblock}{Наївне рішення}
        \footnotesize Список на сервері відкрито $\Rightarrow$ сервер бачить соціальний граф (social graph).
        $\Rightarrow$ Потрібне доведення з нульовим розголошенням.
      \end{alertblock}
    \end{column}
    \begin{column}{0.30\textwidth}
      \centering\photo[width=\linewidth, height=0.74\textheight, keepaspectratio]{signal_group_info.png}
    \end{column}
  \end{columns}
  \footnotetext[1]{Signal Blog, 2019 \cite{SignalBlog2019}}
  \footnotetext[2]{Chase, Perrin, Zaverucha, CCS 2020 \cite{CPZ20}}
\end{frame}

%  5. Проєкт: архітектура і ролі

\begin{frame}{Архітектура і ролі: хто що доводить}
  \begin{columns}[c]
    \begin{column}{0.60\textwidth}
      \centering
      \resizebox{\linewidth}{!}{\begin{tikzpicture}[
    font=\normalsize,
    box/.style={draw=DarkOrchid, fill=Thistle!55, rounded corners=3pt, align=center, minimum height=1.0cm},
    store/.style={draw=DarkOrchid, fill=white, font=\footnotesize, text width=4.2cm, minimum height=0.42cm,
                  align=center},
    flow/.style={-{Stealth[length=2.2mm]}, thick, DarkOrchid},
    lbl/.style={font=\footnotesize, align=center, fill=white, inner sep=1.5pt}
  ]
  \node[box, minimum width=3.1cm] (A) at (0,1.6)  {Аліса (учасниця)\\[-1pt]\footnotesize має $\mathit{GroupMasterKey}$};
  \node[box, minimum width=3.1cm] (B) at (0,-1.6) {Боб (учасник)\\[-1pt]\footnotesize має $\mathit{GroupMasterKey}$};
  \node[box, minimum width=4.8cm, minimum height=4.6cm, fill=Lavender] (S) at (8.2,0) {};
  \node[font=\normalsize\bfseries, anchor=north] at (S.north) {Сервер Signal};
  \node[font=\footnotesize, anchor=north, align=center] at ($(S.north)+(0,-0.45)$)
       {видавець \emph{і} верифікатор\\ має $\mathit{ServerSecretParams}$};
  \node[store] (r1) at ($(S.center)+(0,0.05)$)
       {$\mathsf{Enc}(\mathrm{UID}_A)\;\big|\;\mathsf{Enc}(\mathrm{PK}_A)\;\big|\;$роль};
  \node[store, below=0pt of r1] (r2) {$\mathsf{Enc}(\mathrm{UID}_B)\;\big|\;\mathsf{Enc}(\mathrm{PK}_B)\;\big|\;$роль};
  \node[store, below=0pt of r2] (r3) {$\mathsf{Enc}$(назва), $\mathsf{Enc}$(аватар), \dots};
  \node[font=\footnotesize, above=1pt of r1] {зашифрований стан групи};
  \draw[flow] ($(A.east)+(0,0.25)$) -- node[lbl, above=2pt]
        {\textcircled{\scriptsize 1} щоденний $\mathit{AuthCredential}$\\[-1pt]
        (видано для $\mathrm{UID}_A$)} ($(S.west)+(0,1.85)$);
  \draw[flow] ($(S.west)+(0,1.25)$) -- ($(A.east)+(0,-0.25)$);
  \draw[flow] ($(A.south east)+(-0.2,0)$) -- node[lbl, pos=0.55]
        {\textcircled{\scriptsize 2} $\bigl(\mathsf{Enc}(\mathrm{UID}_A),\ \pi\bigr)$\\[-1pt] анонімний запит}
        ($(S.west)+(0,0.2)$);
  \draw[flow] ($(S.west)+(0,-0.6)$) -- node[lbl, pos=0.5] {\textcircled{\scriptsize 3} зашифровані записи}
        ($(B.north east)+(-0.2,0)$);
  \draw[flow, dashed, Orchid] ($(A.south)+(-0.5,0)$) -- node[lbl, left=1pt] {ключ передано\\ наскрізно (E2E)}
        ($(B.north)+(-0.5,0)$);
\end{tikzpicture}
}
    \end{column}
    \begin{column}{0.38\textwidth}
      \small
      \begin{tabular}{@{}p{0.34\linewidth}p{0.62\linewidth}@{}}
        \toprule
        Доказувач (prover)       & клієнт учасника \\
        Верифікатор (verifier)   & сервер Signal \\
        Видавець (issuer)        & \textbf{той самий} сервер $\Rightarrow$ KVAC\footnotemark[1] \\
        Приховано                & UID, ключ профілю, зв'язок між запитами \\
        Розкрито                 & день $\tau$, шифротекст $E$ \\
        \bottomrule
      \end{tabular}\par\medskip
      \begin{block}{Твердження (statement)}
        \footnotesize "Маю дійсне посвідчення на \emph{якийсь} UID, і $E$ зашифровано \emph{саме цим} UID 
        ключем \emph{цієї} групи"{}
      \end{block}
    \end{column}
  \end{columns}
  \footnotetext[1]{Chase, Perrin, Zaverucha, CCS 2020~\cite{CPZ20}}
\end{frame}

%  6. Протокол: Σ-протокол

\section{Протокол}
\begin{frame}{Один двигун для всіх доведень: $\Sigma$-протокол}
  \begin{columns}[T]
    \begin{column}{0.56\textwidth}
      \centering
      \resizebox{\linewidth}{!}{\begin{tikzpicture}[
    font=\normalsize,
    party/.style={draw=DarkOrchid, fill=Thistle!60, rounded corners=3pt, minimum width=3.0cm, minimum height=0.75cm,
                  font=\normalsize\bfseries, align=center},
    msg/.style={-{Stealth[length=2.2mm]}, thick, DarkOrchid},
    note/.style={font=\small, align=center}
  ]
  \node[party] (P) at (0,0)   {Доказувач $\mathcal{P}$\\[-2pt]\footnotesize\mdseries(prover)};
  \node[party] (V) at (7.4,0) {Верифікатор $\mathcal{V}$\\[-2pt]\footnotesize\mdseries(verifier)};
  \node[note, below=1mm of P] {знає $w\in H$: $\psi(w)=t$};
  \node[note, below=1mm of V] {знає $(\psi, t)$};
  \node[note] at (0,-1.55)   {$r \xleftarrow{\$} H$,\; $a \gets \psi(r)$};
  \node[note] at (7.4,-2.45) {$c \xleftarrow{\$} \mathcal{C}$};
  \node[note] at (0,-3.35)   {$z \gets r + c\cdot w$};
  \node[note] at (7.4,-4.25) {прийняти, якщо $\psi(z) = a \cdot t^{\,c}$};
  \draw[msg] (1.7,-1.55) -- node[above, font=\small] {зобов'язання $a$ (commitment)} (5.7,-1.55);
  \draw[msg] (5.7,-2.45) -- node[above, font=\small] {виклик $c$ (challenge)} (1.7,-2.45);
  \draw[msg] (1.7,-3.35) -- node[above, font=\small] {відповідь $z$ (response)} (5.7,-3.35);
  \node[draw=Orchid, fill=Lavender, rounded corners=3pt, align=left, font=\small, anchor=north] at (3.7,-4.8)
    {\textbf{Фіат--Шамір (як у \texttt{poksho}):}
     $c \gets H\bigl(\text{мітка}\,\|\,\text{твердження}\,\|\,\text{точки}\,\|\,a\,\|\,\text{повідомлення}\bigr)$\\[1pt]
     доведення $\pi=(c,z)$; $\mathcal{V}$ відновлює $a=\psi(z)\cdot t^{-c}$ і хешує};
\end{tikzpicture}
}
    \end{column}
    \begin{column}{0.42\textwidth}
      \begin{block}{Прообраз гомоморфізму}
        Довести знання\footnotemark[1] $w$ такого, що $\psi(w)=t$
      \end{block}
      \begin{itemize}\small
        \item повнота, спец.\ обґрунтованість, HVZK
        \item AND-композиція: спільний виклик
        \item Фіат--Шамір\footnotemark[2]: $c = H(\dots)$
      \end{itemize}
      \begin{alertblock}{У \texttt{zkgroup}}
        \small кожне доведення таке: без схем (circuits) і довіреного налаштування (trusted setup)
      \end{alertblock}
    \end{column}
  \end{columns}
  \footnotetext[1]{Лекція 3, "$\Sigma$-протоколи"{}~\cite{SigmaProtocols}}
  \footnotetext[2]{Fiat, Shamir, CRYPTO 1986~\cite{FS86}}
\end{frame}

%  7. Протокол: алгебраїчний MAC

\begin{frame}{Алгебраїчний MAC (algebraic MAC)}
  \begin{columns}[T]
    \begin{column}{0.52\textwidth}
      \begin{block}{CMZ, 2014\footnotemark[1]: $\mathsf{MAC}_{\mathsf{GGM}}$}
        ключ $(x_0,x_1)$; \ тег на $m\in\Zq$: \ $\sigma = (U,\ U^{x_0+x_1 m})$
      \end{block}
      \begin{block}{Signal (CPZ, 2020)\footnotemark[2]: атрибути це \emph{елементи групи}}
        $V = W \cdot U^{x_0+x_1 t}\cdot \prod_i M_i^{\,y_i}$
      \end{block}
      \begin{alertblock}{Навіщо}
        \small Дешеве верифіковане шифрування (verifiable encryption):\\ \textbf{64 байти} замість $>600$
      \end{alertblock}
    \end{column}
    \begin{column}{0.45\textwidth}
      \begin{exampleblock}{Іграшковий приклад: $q=11$ у $\mathbb{Z}_{23}^{\times}$, $g=2$}
        \small
        ключ $(3,5)$, \ $m=2$: \ $x_0+x_1m = 13 \equiv 2$\\[2pt]
        тег: $(U,U') = (4,\ 4^2) = (4,16)$\\[4pt]
        перерандомізація (re-randomisation), $a=4$:\\
        $(4^4,\ 16^4) = (3,\ 9)$, \quad $3^2 = 9$ \checkmark\\[4pt]
        \textbf{новий тег на те саме $m$} -- видавець не пов'яже його з $(4,16)$
      \end{exampleblock}
    \end{column}
  \end{columns}
  \footnotetext[1]{Chase, Meiklejohn, Zaverucha, CCS 2014~\cite{CMZ14}}
  \footnotetext[2]{Chase, Perrin, Zaverucha, CCS 2020~\cite{CPZ20}}
\end{frame}

%  8. Протокол: видача і пред'явлення

\begin{frame}{Видача і пред'явлення (issue \& present)}
  \begin{columns}[T]
    \begin{column}{0.54\textwidth}
      \centering
      \resizebox{\linewidth}{!}{\begin{tikzpicture}[
    font=\normalsize,
    party/.style={draw=DarkOrchid, fill=Thistle!60, rounded corners=3pt, minimum width=3.0cm, minimum height=0.75cm,
                  font=\normalsize\bfseries, align=center},
    msg/.style={-{Stealth[length=2.2mm]}, thick, DarkOrchid},
    lab/.style={above, font=\footnotesize, fill=white, inner sep=1pt},
    phase/.style={draw=Orchid, fill=Lavender, rounded corners=2pt, font=\footnotesize\bfseries, inner sep=2pt},
    loc/.style={font=\footnotesize, align=left}
  ]
  \node[party] (U) at (0,0)    {Користувач\\[-2pt]\footnotesize(клієнт)};
  \node[party] (S) at (9.0,0)  {Сервер\\[-2pt]\footnotesize видавець $=$ верифікатор};
  \draw[densely dotted, Plum] (U.south) -- (0,-8.7);
  \draw[densely dotted, Plum] (S.south) -- (9.0,-8.7);
  \node[phase] at (4.5,-0.95) {Видача (issue): з автентифікацією, раз на добу};
  \draw[msg] (0.1,-1.75) -- node[lab] {ідентифікатор $\mathrm{UID}$ (звичайний вхід)} (8.9,-1.75);
  \node[loc, anchor=east] at (8.85,-2.55) {$t\xleftarrow{\$}\mathbb{Z}_q,\ U\xleftarrow{\$}\mathbb{G}$,\\
         $V = W\,U^{x_0+x_1t}\prod_i M_i^{y_i}$};
  \draw[msg] (8.9,-3.45) -- node[lab] {$(t,U,V)$ і доведення $\pi_I$ коректної видачі} (0.1,-3.45);
  \node[loc, anchor=west] at (0.15,-3.95) {перевірити $\pi_I$ за публічними $(C_W, I)$};
  \node[phase] at (4.5,-4.7) {Пред'явлення (present): анонімно, скільки завгодно разів};
  \node[loc, anchor=west] at (0.15,-5.5) {$z\xleftarrow{\$}\mathbb{Z}_q$:\ $C_{x_0}=G_{x_0}^zU,\ C_{x_1}=G_{x_1}^zU^t$,\\
         $C_{y_i}=G_{y_i}^zM_i,\ C_V=G_V^zV$;\ шифротекст $E$};
  \draw[msg] (0.1,-6.45) -- node[lab] {$\bigl(C_{x_0},C_{x_1},\{C_{y_i}\},C_V,E,\pi\bigr)$} (8.9,-6.45);
  \node[loc, anchor=east] at (8.85,-7.25) {$Z \gets C_V\big/\bigl(W\,C_{x_0}^{x_0}C_{x_1}^{x_1}\prod_i C_{y_i}^{y_i}\bigr)$,\\
         прийняти, якщо $\pi$ перевіряється з цим $Z$};
  \node[loc, anchor=west, DarkOrchid, text width=7.6cm] at (0.15,-8.2)
       {нове $z$ щоразу $\Rightarrow$ пред'явлення не зв'язуються між собою і з видачею (unlinkability)};
\end{tikzpicture}
}
    \end{column}
    \begin{column}{0.44\textwidth}
      \begin{block}{Ключова рівність}
        \centering $Z = I^{z}$ \quad $\Longleftrightarrow$ \quad MAC дійсний
      \end{block}
      \small Доведення\footnotemark[1] $\pi_A$ -- одна AND-композиція:\par\smallskip
      \footnotesize
      \begin{tabular}{@{}ll@{}}
        $Z = I^{z}$                                   & дійсне посвідчення \\
        $C_{x_1} = C_{x_0}^{\,t}G_{x_0}^{z_0}G_{x_1}^{z}$ & узгоджене $t$ \\
        $A = G_{a_1}^{a_1}G_{a_2}^{a_2}$              & ключ \emph{цієї} групи \\
        $E_1 = C_{y_1}^{\,a_1}G_{y_1}^{z_1}$          & $E_1 = M_1^{a_1}$ \\
        $C_{y_2}/E_2 = G_{y_2}^{z}/E_1^{\,a_2}$        & $E_2 = E_1^{a_2}M_2$ \\
      \end{tabular}\par\smallskip
      \small усі рівняння \textbf{лінійні} в секретних показниках
    \end{column}
  \end{columns}
  \footnotetext[1]{Chase, Perrin, Zaverucha, CCS 2020, §5~\cite{CPZ20}}
\end{frame}

%  9. Протокол: реалізація

\begin{frame}{Реалізація: бібліотека \texttt{poksho}}
  \begin{columns}[T]
    \begin{column}{0.40\textwidth}
      \begin{itemize}\footnotesize
        \item \textbf{Сильний Фіат--Шамір} (strong FS): гешується опис твердження $+$ усі точки\footnotemark[1]
        \item \textbf{Синтетичні нонси} (synthetic nonces): $r$ з випадковості \emph{і} гешу секрету
        \item \textbf{Fault attacks}; доведення $(c,\vec z)$: $32(1+k)$ Б
        \item \textbf{ristretto255}\footnotemark[2]: простий порядок, без кофактора
      \end{itemize}
      \begin{alertblock}{Слабкий FS ламав системи}
        \footnotesize $c = H(a)$ $\Rightarrow$ твердження підбирають \emph{після} виклику: Helios 2012\footnotemark[3],
        Frozen Heart 2022\footnotemark[4]
      \end{alertblock}
    \end{column}
    \begin{column}{0.58\textwidth}
      \centering\photo[width=\linewidth, height=0.80\textheight, keepaspectratio]{poksho_absorb.png}
    \end{column}
  \end{columns}
  \footnotetext[1]{libsignal, \texttt{poksho/src/statement.rs} \cite{libsignal}}
  \footnotetext[2]{RFC 9496 \cite{RFC9496}}
  \footnotetext[3]{Bernhard, Pereira, Warinschi, ASIACRYPT 2012 \cite{BPW12}}
  \footnotetext[4]{Trail of Bits, 2022 \cite{FrozenHeart22}}
\end{frame}

%  10. Аналіз: припущення і ризики

\section{Аналіз}
\begin{frame}{Припущення і ризики}
  \begin{columns}[T]
    \begin{column}{0.32\textwidth}
      \begin{block}{Припущення\footnotemark[1]}
        \small
        \begin{itemize}
          \item DDH у $\G$
          \item MAC$_\mathsf{GGM}$: доведено лише в моделі загальної групи (GGM)
          \item Модель випадкового оракула (ROM)
        \end{itemize}
      \end{block}
    \end{column}
    \begin{column}{0.32\textwidth}
      \begin{block}{Властивості}
        \small
        \begin{itemize}
          \item Ключ сервера $=$ право "друкувати"{} посвідчення
          \item Детерміноване шифрування: сервер бачить рівність записів
          \item Cкладна not black-box композиція
        \end{itemize}
      \end{block}
    \end{column}
    \begin{column}{0.32\textwidth}
      \begin{alertblock}{Квантова загроза}
        \small
        \begin{itemize}
          \item Усе на стійкості дискретного логарифму $\Rightarrow$ алгоритм Шора
          \item HNDL: "harvest now -- decrypt later"{}
        \end{itemize}
      \end{alertblock}
    \end{column}
  \end{columns}
  \vfill
  \centering\small Публічних звітів про вразливості саме в \texttt{zkgroup} в інтернеті немає, але це не гарантія 
  безпеки!
  \footnotetext[1]{Chase, Perrin, Zaverucha, CCS 2020~\cite{CPZ20}}
\end{frame}

%  11. Аналіз: продуктивність

\begin{frame}{Продуктивність: власні виміри}
  \begin{columns}[T]
    \begin{column}{0.60\textwidth}
      \scriptsize\setlength{\tabcolsep}{3.5pt}
      \begin{tabular}{@{}lrrr@{}}
        \toprule
        Об'єкт                   & \shortstack[r]{байти                      \\2020\footnotemark[1] $\to$ 2026} & \shortstack[r]{створення,\\мс}
                                 & \shortstack[r]{перевірка,                 \\мс} \\
        \midrule
        UidCiphertext            & 64 $\to$ 65               & 0.085 & 0.158 \\
        AuthCredential: response & 361 $\to$ 425             & 1.15  & 0.51  \\
        AuthCredential: present  & 493 $\to$ 633             & 1.85  & 0.95  \\
        ProfileKeyCred.: present & 713 $\to$ 721             & 1.97  & 1.12  \\
        \bottomrule
      \end{tabular}\par\smallskip
      Intel Core i7-1065G7, Windows; \texttt{cargo bench} (criterion)\footnotemark[2]\par\smallskip
      \centering\photo[width=\linewidth, height=0.42\textheight, keepaspectratio]{cargo_bench_terminal.png}
    \end{column}
    \begin{column}{0.38\textwidth}
      \begin{block}{Чому байтів більше, ніж у статті?}
        \footnotesize
        \begin{itemize}
          \item \textbf{+1 байт}: версія формату
          \item \textbf{493 $\to$ 633}: Посвідчення засвідчує \emph{два} ідентифікатори (ACI $+$ PNI)
          \item Розміри не залежать від процесора
        \end{itemize}
      \end{block}
      \begin{alertblock}{Підсумок}
        \footnotesize усе $\approx$ 1--2 мс, $<1$ КБ; група зі 100 учасників $\approx$ 0.07 с
      \end{alertblock}
    \end{column}
  \end{columns}
  \footnotetext[1]{Chase, Perrin, Zaverucha, CCS 2020, Table 1 \cite{CPZ20}}
  \footnotetext[2]{libsignal, коміт 8509425 \cite{libsignal}}
\end{frame}

%  12. Аналіз: альтернатива BBS

\begin{frame}{Чи можна замінити на BBS?}
  \begin{columns}[T]
    \begin{column}{0.54\textwidth}
      \centering
      \resizebox{\linewidth}{!}{\begin{tikzpicture}
  \begin{axis}[
      ybar, bar width=13pt, width=9.5cm, height=5.6cm,
      ymin=0, ymax=15.5, ylabel={час, мс}, ylabel style={font=\small},
      symbolic x coords={Видача/підпис, Створення доведення, Перевірка},
      xtick=data, xticklabel style={font=\small, align=center, text width=2.4cm},
      yticklabel style={font=\footnotesize},
      enlarge x limits=0.22,
      legend style={at={(0.5,1.03)}, anchor=south, font=\footnotesize, draw=Plum, legend columns=1},
      nodes near coords, nodes near coords style={font=\scriptsize, /pgf/number format/.cd, fixed, precision=2},
      ymajorgrids, grid style={Thistle!70}
    ]
    \addplot[draw=DarkOrchid, fill=Orchid!70] coordinates
      {(Видача/підпис,1.15) (Створення доведення,1.85) (Перевірка,0.95)};
    \addplot[draw=Plum!80!black, fill=Thistle] coordinates
      {(Видача/підпис,7.04) (Створення доведення,13.05) (Перевірка,12.42)};
    \legend{{zkgroup AuthCredential (ristretto255, KVAC)},{BBS, $L=4$, $U=3$ (BLS12-381, спарювання)}}
  \end{axis}
\end{tikzpicture}
}\par
      \scriptsize та сама машина; BBS -- \texttt{zkryptium}\footnotemark[1]; показано \emph{порядок}, не точне
      співвідношення
    \end{column}
    \begin{column}{0.44\textwidth}
      \small
      \textbf{BBS}\footnotemark[2]: підпис 80 Б, доведення $272 + 32U$ Б, $U$ -- undisclosed messages, стандарт 
      IETF у розробці\par\medskip
      \textbf{Ціна заміни:}
      \begin{enumerate}
        \item Спарювання (pairings): на порядок повільніше (перевірено реалізацією)
        \item Потреба другої кривої (BLS12-381)
        \item \alert{Повідомлення BBS -- скаляри} $\Rightarrow$ ламається дешеве шифрування точок
      \end{enumerate}
      \textbf{Виграш:} Публічна верифікація -- Signal не потрібна\par\medskip
      Як компроміс: BBS-MAC / $\mu$BBS\footnotemark[3] -- без спарювань
    \end{column}
  \end{columns}
  \footnotetext[1]{zkryptium 0.8.1~\cite{zkryptium}}
  \footnotetext[2]{IRTF CFRG, draft-irtf-cfrg-bbs-signatures-12~\cite{BBSdraft}}
  \footnotetext[3]{Orr{\`u}, ePrint 2024/1552 (CCS 2025)~\cite{Orru24}}
\end{frame}

%  13. Висновки

\section{Висновки}
\begin{frame}{Висновки}
  \begin{enumerate}
    \item Модель довіри визначає криптографію: видавець $=$ верифікатор $\Rightarrow$ \textbf{KVAC}, а не підписи
    \item Внесок Signal: MAC на \textbf{елементах групи} $\Rightarrow$ дешеве верифіковане шифрування
    \item Виміри: $\approx$ 1-2 мс, $<1$ КБ на операцію
    \item Слабкі місця: GGM, складність композиції, \alert{відсутність постквантової стійкості}
    \item BBS тут -- крок назад; реальна еволюція -- нові KVAC або постквантовий редизайн
  \end{enumerate}
  \note{
    Підсумую доповідь п'ятьма тезами. Перше: криптографію визначає модель довіри -- оскільки сервер і видає, 
    і перевіряє, Signal обрав верифікацію за ключем. Друге: головний винахід -- MAC, у якого атрибути є точками 
    групи; це зробило шифрування учасників компактним. Третє: це швидко -- мілісекунди і сотні байтів. 
    Четверте: слабкі місця -- теоретична модель загальної групи, складність і, головне, квантова загроза. 
    І п'яте: стандартний BBS тут не допоможе; майбутнє за новими схемами з верифікацією за ключем або 
    постквантова перебудова всієї системи.}
\end{frame}

% 14. Питання для самоперевірки

\begin{frame}{Питання для самоперевірки}
  \small
  \begin{enumerate}
    \item Чому Signal не потрібна публічна верифікованість?\\
          \rotatebox[origin=c]{180}{{\color{DarkOrchid}(Видавець і верифікатор -- той самий сервер, тож досить KVAC на алгебраїчному MAC.)}}
    \item Що означає рівність $Z = I^z$ при пред'явленні?\\
          \rotatebox[origin=c]{180}{{\color{DarkOrchid}(Клієнт має дійсний MAC сервера на приховані атрибути; доводиться як Шнорр.)}}
    \item Чому атрибути-точки кращі за атрибути-скаляри?\\
          \rotatebox[origin=c]{180}{{\color{DarkOrchid}(Дають дешеве верифіковане шифрування UID: 64 байти замість $>600$.)}}
    \item Чим небезпечний слабкий Фіат--Шамір $c = H(a)$?\\
          \rotatebox[origin=c]{180}{{\color{DarkOrchid}(Твердження можна підібрати після виклику $\Rightarrow$ "доведення"{} хибних тверджень.)}}
    \item Що буде, якщо доказувач двічі використає той самий нонс $r$?\\
          \rotatebox[origin=c]{180}{{\color{DarkOrchid}(Спеціальна обґрунтованість дає секрет: $x = (z - z')/(c - c')$.)}}
    \item Чому \texttt{zkgroup} вразливий у постквантову епоху?\\
          \rotatebox[origin=c]{180}{{\color{DarkOrchid}(Усе на дискретному логарифмі; загроза HNDL для збережених шифротекстів груп.)}}
  \end{enumerate}
\end{frame}

\begin{frame}[allowframebreaks]{Джерела}
  \renewcommand*{\bibfont}{\tiny}
  \renewcommand*{\mkbibacro}[1]{#1}
  \DeclareFieldFormat{url}{\href{#1}{\mkbibacro{(URL)}}}
  \printbibliography[heading=none]
\end{frame}

\end{document}
